% Standalone mathematical contribution; no claim of a completed Lean formalization.
\section{The PA-internal converse for the PA196 command calculus}
\label{sec:internal-converse}

The published endpoint of PA196 is
\[
  \mathrm{PA}\vdash \neg H_{196}\longrightarrow\operatorname{Con}(\mathrm{PA}).
\]
Its separate external nonhalting theorem uses a concrete model of hereditarily
finite lists.  That external theorem, by itself, does not establish the
opposite implication inside PA.  The argument below supplies the opposite
implication by a syntactic interpretation.  It is an ordinary mathematical
proof and a specification for an additional formalization; it is not a claim
that the existing Lean declarations already contain this converse.

\subsection{Conventions and the translation}

Write $\mathcal C$ for the numeric calculus of
\texttt{PA196/Logic/Basic.lean}. Its rules are modus ponens and theorem
generalization; its axioms are the six displayed heads, the restricted
propositional head, and the guarded gate. A derivation is a finite tree
whose nodes record the rule and its numerical arguments. Let
$\operatorname{Prf}_{\mathcal C}(p,n)$ be the primitive recursive predicate
saying that $p$ is such a tree with conclusion code $n$.

The translation $A\mapsto A^*$ decodes numeric formulas using the actual
transposed universal constructor. It replaces membership by binary-bit
membership, equality by arithmetic equality, implication by implication,
and a quantifier with label $a$ by $\forall x_a$. Precisely,
\[
  \mathsf M(a,b)^*=\operatorname{Bit}(x_a,x_b),\qquad
  \mathsf E(a,b)^*=(x_a=x_b),\qquad
  \mathsf U(a,B)^*=\forall x_a\,B^*.
\]
Here $\operatorname{Bit}(s,i)$ means that the binary digit of $s$ at
position $i$ is one. Code zero is decoded as $\top$, for example $0=0$.
The numeric identity $\mathsf I(0,0)=0$ is the sole degenerate implication
case: the translation proves $0^*\leftrightarrow(0^*\to0^*)$, rather than
pretending that their parse trees are identical. For every other implication
code, decoding gives the indicated implication literally. Consequently the
translation of an MP step is either ordinary MP, or the fixed proof of
$\top$ in this exceptional case. Code one is
$\mathsf M(0,0)$ and hence translates to
$\operatorname{Bit}(x_0,x_0)$.

We may carry out the proof in the conservative definitional extension
$\mathrm{PA}^{\mathrm{pr}}$ with symbols for the fixed primitive recursive
functions needed below, including exponentiation and binary digits.
Their totality is provable in PA, and the usual elimination of those
definitions is a primitive recursive proof transformation. No uniform
evaluation of arbitrary primitive recursive programs is being postulated.
Only finitely many fixed numerical functions occur in these definitions.

\subsection{The finite-comprehension proof template}

For every arithmetic formula $\varphi(i,\vec z)$, with $s,n$ chosen fresh,
PA proves the corresponding instance
\begin{equation}\label{eq:finite-comprehension-template}
 \forall n\;\exists s<2^n\;\forall i<n\,
 \bigl(\operatorname{Bit}(s,i)\leftrightarrow\varphi(i,\vec z)\bigr).
\end{equation}
This uses full induction, so no restriction is placed on the quantifier
complexity of $\varphi$. At $n=0$, choose $s=0$. In the induction step,
choose $s'=s+2^n$ if $\varphi(n,\vec z)$, and $s'=s$ otherwise.
The bound $s<2^n$ ensures that adding $2^n$ changes exactly the new bit;
the familiar bit identities and $s'<2^{n+1}$ complete the step. Classical
excluded middle for $\varphi(n,\vec z)$ is available in PA.

This is a syntactic template, not a uniform truth definition. Given the
code of $\varphi$, a primitive recursive algorithm writes the displayed
induction instance and the fixed proof around it. Fresh variable selection,
capture-avoiding substitution, the induction-axiom recognizer, and proof
concatenation are primitive recursive. PA proves the algorithm's output is
a proof of the corresponding instance by checking these syntactic
operations. The induction instance changes with $\varphi$; no assertion
that PA proves its own soundness occurs.

PA also proves
\begin{equation}\label{eq:no-self-bit}
 \forall a\;\neg\operatorname{Bit}(a,a),
 \qquad
 \operatorname{Bit}(a,i)\longrightarrow i<a.
\end{equation}
For the second assertion, a set bit implies $2^i\le a$, while
$i<2^i$ is provable by induction. The first follows by setting $i=a$.

\subsection{The gate, including the coincident-label case}

Fix numerical labels $a,b,c$ with $c\ne a,b$. Put
$\Phi=(\operatorname{decode}(c))^*$. Every variable label appearing
anywhere in $\operatorname{decode}(c)$ is strictly below $c$.
The reason is that a nonzero constructor's Cantor components are smaller
than the constructor code, in either order; induction on the decoder
proves the assertion. In particular $x_c$ does not occur in $\Phi$, even
bound. This syntactic fact is primitive recursively checkable and
PA-provable. It is stronger than the required freshness condition.

First suppose $a\ne b$. Under fixed values $A=x_a$ and $B=x_b$, apply
\eqref{eq:finite-comprehension-template} at $n=A+B+1$ to
\[
  \theta(i)=
  \bigl(\operatorname{Bit}(A,i)\lor i=B\bigr)
  \land \Phi[i/x_b].
\]
The resulting number $s$ satisfies
\[
 \forall i\;
   \bigl(\operatorname{Bit}(A,i)\to\Phi[i/x_b]
                 \to\operatorname{Bit}(s,i)\bigr)
   \land\bigl(\operatorname{Bit}(s,i)\to\Phi[i/x_b]\bigr),
 \qquad
 \Phi\to\operatorname{Bit}(s,B).
\]
For $i\ge n$, both the first antecedent and $\operatorname{Bit}(s,i)$
are false, by the bound on $s$ and \eqref{eq:no-self-bit}; hence the
unbounded universal statement follows from the bounded comprehension.

If $a=b$, the occurrence of $x_a$ in the gate's universally quantified
membership antecedent is bound together with $x_b$. That antecedent is
then $\operatorname{Bit}(i,i)$, always false. In this case use
$n=B+1$ and $\theta(i)=(i=B)\land\Phi[i/x_b]$ instead. The same required
separation statement and distinguished-point implication follow. This
case distinction is necessary: silently treating $x_a$ as a fixed
parameter when $a=b$ would be a variable-capture error.

These two cases give precisely the translation of
\[
   \exists x_c\bigl(S(a,b,c)\land P(a,b,c)\bigr),
\]
where $P=c\to\mathsf M(c,b)$ and
\[
  S=\mathsf U\bigl(b,
      (\mathsf M(a,b)\to P)\land(\mathsf M(c,b)\to c)\bigr).
\]
Because $c\ne a,b$, the equality- and membership-vacuity clauses
$\Delta$ and $V$ are ordinary quantifier-vacuity tautologies after
translation. If $Q=\mathsf U(c,S\to\neg P)$, the displayed witness
refutes $Q$. From $\Delta$, $V$, and $\neg Q$, propositional logic gives
\[
   \neg\bigl(\Delta\to(V\to Q)\bigr),
\]
which is the exact PA196 gate. When a guard fails, the command returns
code zero, already covered by the proof of $\top$.

After using the fixed PA proof that $\operatorname{Bit}(x_0,x_0)$ is false,
all six other heads translate into first-order logical or equality
theorems. This replacement matters for heads containing encoded negation or
existential quantification: their translations are PA theorems, not pure
logical validities with an arbitrary predicate in place of falsity.
The restricted propositional head translates into a classical
tautology; its numerical restriction is irrelevant to soundness.
The proof templates also cover the exceptional numeric implication code
by the fixed equivalence noted above.

\subsection{A proof compiler and its internal consequence}

There is a primitive recursive function $T$ for which PA proves
\begin{equation}\label{eq:proof-compiler-soundness}
 \operatorname{Prf}_{\mathcal C}(p,1)
 \longrightarrow
 \operatorname{Prf}_{\mathrm{PA}}(T(p),\ulcorner\bot\urcorner).
\end{equation}
To construct it, traverse the derivation tree. Translate each axiom by
the templates above, each MP node by MP (with the exceptional truth case
handled explicitly), and each theorem-generalization node by arithmetic
generalization. At the root append the fixed proof of
$\neg\operatorname{Bit}(x_0,x_0)$ from \eqref{eq:no-self-bit}, followed
by the contradiction rule of the selected PA presentation. Eliminate the
finitely many definitional extensions. Recursion on a finite tree and
substitution into fixed proof templates are primitive recursive;
duplication of previously constructed proofs may enlarge certificates
but does not change this fact. PA proves correctness by induction on the
tree code, using its ordinary syntactic proof-checking predicates.

Thus PA proves
\begin{equation}\label{eq:con-no-calculus-refutation}
 \operatorname{Con}(\mathrm{PA})\longrightarrow
 \forall p\;\neg\operatorname{Prf}_{\mathcal C}(p,1).
\end{equation}
The existing HF interpretation is unnecessary for this implication.

\subsection{From proof syntax to the literal machine}

The sound-store invariant needed by the command interpreter is purely
syntactic: its live formulas have $\mathcal C$ derivations and every
element of the stack has one too. A zero tail pops to zero, and zero has
a fixed derivation. An AX step attaches the appropriate axiom proofs
and the theorem-generalization proof. An MP step either attaches an
MP proof or returns zero. A Shift step halts only if $t+h=1$, which
forces one of the two live formulas to be code one.
Consequently PA proves, by induction on a coded finite command
execution, that a halting execution from the zero store gives a
$\mathcal C$ refutation. This proof uses the existence of finite proof
codes, not their truth in a model.

To transport this assertion to a specific literal table, one must also
internalize its startup and instruction contracts. The relevant
contracts are finite, positive-length simulations on clean unary
register frames. Their variable parts are the fixed tape-engine
loops, proved by arithmetic induction; their program-reading parts
are finite symbolic traces. All register operations, path encodings,
and trace concatenations are primitive recursive, and PA proves
their totality. They therefore hold also for nonstandard coded frames
and traces in arbitrary models of PA. Determinism and positive length
exclude an intermediate halt before a prescribed return frame. More
explicitly, if a literal run halts within $t$ steps while none of the
first $t+1$ source instructions halts, the total source-step function
and the positive instruction contracts produce a literal run lasting
at least $t+1$ steps from that frame. Determinism contradicts the
claimed earlier halt. This is a bounded induction argument available
inside PA; it does not invoke the truth of an unbounded nonhalting
assertion. The identical argument applies to command boundaries using
the positive, total arithmetic macro contracts. This
gives, within PA, equivalence of blank-tape halting and source-program
halting, and hence the literal implication
\[
 \mathrm{PA}\vdash H_{196}\longrightarrow
              \neg\operatorname{Con}(\mathrm{PA}).
\]
Together with the published reverse implication this proves
\[
 \mathrm{PA}\vdash
 \neg H_{196}\ \longleftrightarrow\ \operatorname{Con}(\mathrm{PA}).
\]
For a newly reduced table the same conclusion follows once its finite
symbolic instruction contracts and startup preserve this source
program. This mathematical proof should be distinguished from the
stronger artifact claim of a completed Lean formalization of the
new internal converse: no such completed formalization is supplied
by the upstream PA196 endpoint.

\subsection{What transfers from ConciseIPRA}

ConciseIPRA is useful here as a design principle: replace primitive
logical and equality heads by derivations, exploit a deliberately
chosen representation, and count the implementation of the resulting
checker rather than surface axiom symbols. PA196's removal of the
primitive seventh equality head is exactly this kind of reduction.
The mathematical target cannot, however, simply be replaced by the
ConciseIPRA calculus. Its stated equivalence is with intuitionistic
primitive recursive arithmetic, whose induction is limited to
equations after the appropriate translation. Our finite-comprehension
template deliberately uses induction for arbitrary arithmetic
formulas. A machine searching only for an iPRA refutation would
establish the consistency of that theory; the required PA-internal
equivalence with $\operatorname{Con}(\mathrm{PA})$ would require an
additional theorem and does not follow from ConciseIPRA's
bi-interpretation theorem.
